\section{Positions Synopsis}

The table below distinguishes the philosophical, scientific, and Catholic doctrinal positions compared across the preceding argument, separating legitimate empirical discipline from illegitimate epistemic exclusion.

\begin{threestable}{Position}{Core Claim}{Distinction and Assessment}
Instrumental exclusivism & Only instrument-detectable entities exist, or only empirically verifiable propositions qualify as genuine knowledge or cognitively meaningful truth. & Target of refutation: presupposes non-metrical logic, math, norms, testimony, and rational judgment; fails under substantive, stipulative, and metalinguistic formulations. \\
\addlinespace[0.35em]
Philosophy-of-science instrumentalism and constructive empiricism & Classical instrumentalism treats unobservable theoretical posits (fields, quarks, wavefunctions) as predictive mathematical tools (e.g., Mach; instrumentalist readings of Duhem); constructive empiricism construes theories literally while limiting acceptance to empirical adequacy (van Fraassen). & Distinct from instrumental exclusivism; both express epistemic modesty regarding physical models rather than denying realities beyond physics. \\
\addlinespace[0.35em]
Ordinary empiricism & Human knowledge begins in sensible experience, and factual claims about the physical world must answer to observation and calibrated testing. & Affirmed by classical realism (Aristotle, Aquinas) and modern science alike; becomes exclusivist only when it denies that intellect rises beyond sense. \\
\addlinespace[0.35em]
Methodological naturalism & Empirical natural sciences restrict their own discipline-specific explanations to secondary physical causes and repeatable measurable regularities. & Lawful craft boundary recognized by Vatican I (\emph{Dei Filius}, ch.\ 4) and Vatican II (\emph{Gaudium et spes}, no.\ 36); becomes illegitimate only when promoted to metaphysical ontology. \\
\addlinespace[0.35em]
Classical logical positivism (verificationism) & Every cognitively meaningful proposition is either an analytic tautology or empirically verifiable (Vienna Circle; A.\ J.\ Ayer, 1936 and 1946). & Collapsed internally: strong verifiability outlawed universal scientific laws; weak verifiability admitted arbitrary statements (Church, 1949; Hempel, 1950); analytic/synthetic pillar fell to Quine (1951). \\
\addlinespace[0.35em]
Popperian falsificationism & Empirical scientific theories are demarcated by exposing themselves to potential observational falsification (Karl Popper, 1934, 1959, 1963). & A criterion of scientific demarcation, never of meaning or reality; Popper explicitly defended the meaningfulness and rational arguability of metaphysics. \\
\addlinespace[0.35em]
Quinean holism & Statements face sensory experience corporately as an interconnected web of belief rather than sentence by sentence; no sharp analytic/synthetic wall holds (W.\ V.\ O.\ Quine, 1951). & Dismantled classical verificationist reductionism while illustrating that logic, ontology, and empirical inquiry cannot be cleanly segregated by a meter. \\
\addlinespace[0.35em]
Thomistic epistemic realism and \emph{demonstratio quia} & Embodied intellect begins in sense, abstracts through physical, mathematical, and metaphysical orders, and reasons \emph{a posteriori} from sensible effects to their transcendent First Cause (\emph{ST} I, q.\ 2, a.\ 2; q.\ 12, a.\ 12; qq.\ 79, 84--85). & Establishes \emph{that} God exists (\emph{an sit}) and what belongs to him as First Cause exceeding all effects, without pretending to comprehend his essence (\emph{quid sit}) or bypassing empirical reality. \\
\addlinespace[0.35em]
Catholic magisterial teaching on reason, faith, and science & God can be known with certainty by natural reason from created things; faith and reason form a twofold harmonious order; each science enjoys just liberty in its own sphere (\emph{Dei Filius}, chs.\ 1, 2, 4; \emph{Aeterni Patris}; \emph{Gaudium et spes} 36, 57; \emph{CCC} 31--39; \emph{Fides et ratio} 55, 67, 83, 88). & Solemnly defines reason's natural capacity without canonizing a single philosophical school; rejects both scientism (shrinking reason to positive science) and fideism (severing faith from reason). \\
\addlinespace[0.35em]
Modernist agnosticism and fideism & Reason is confined solely to appearances (\emph{phenomena}); religion flees from intellect into subjective feeling and mutable symbols (diagnosed in St.\ Pius X, \emph{Pascendi Dominici Gregis}, 1907). & The mirror image of scientistic exclusivism: surrendering objective reality to phenomenalism destroys rational credibility and reduces faith to immanent emotion. \\
\end{threestable}
